% Drawing 152/01-9 -- folder 152, batch 01, archive page 1 (PDF page 2).
% Transcription: à gauche, sous le cercle aux faisceaux, au crayon : deux arbres de petites flèches mises bout à bout (le premier monte d'un point marqué par quatre flèches, avec trois flèches latérales vers la droite ; le second, issu de la première flèche latérale, monte par trois flèches avec deux départs latéraux, et une flèche lui arrive par-dessous d'un second point marqué) ; à leur droite, deux points marqués joints par deux traits fins et par deux chaînes de trois et quatre flèches, l'une au-dessus, l'autre au-dessous.
% Kind: pencil trees of small arrows; two points joined by two chains of arrows. Method 2 (schematic redraw from the facsimile). Status: ok.
% Notes: First tree: the marked point at the bottom left, four arrows end to end going up from it, three side arrows to the right (the top one curved, the bottom one leaving the marked point) and a faint curve beside the second arrow. Second tree: it starts where the bottom side arrow ends; three arrows going up, a side arrow to the right, a faint stroke to the upper right from the second node; under it an arrow coming up from a second marked point, with a faint curve beside it. Right: two marked points joined by two thin straight lines and by two chains of arrows, three above and four below. 21 arrows in all, each drawn with its gap as on the page (the arrows do not touch). Overlaid on the 400 dpi crop. The second marked point is also the top of the three-branched tree of 01-10. Redrawn 2026-11-12 (issue 128): to check again.
% Coordinates: 1 unit = 1 mm = 1/100 of the crop width (bbox in page fractions: [0.14, 0.175, 0.4, 0.345]).
% Redrawn by Claude (claude-fable-5-2), issue #128, 2026-11-12; not checked by a human; not compiled (no TeX in the session).
% Rights: reproduces part of the work -- behind RENDER_TRANSCRIPTS, never published by itself.
\documentclass[tikz,border=3pt]{standalone}
\usepackage{amsmath,amssymb}
\usetikzlibrary{patterns,arrows.meta}
\tikzset{pencil/.style={black!45}, struck/.style={black!35}, mpt/.style={circle,draw,inner sep=1.3pt}, pt/.style={circle,fill,inner sep=0.9pt}, opt/.style={circle,draw,fill=white,inner sep=1.1pt}}
\begin{document}
\begin{tikzpicture}[x=1mm,y=1mm,line cap=round,line join=round,font=\small]
\node[pt] at (13,-69.5) {};
\draw[-{Stealth[length=1.8mm]}, pencil] (13,-69.5) -- (10.7,-55);
\draw[-{Stealth[length=1.8mm]}, pencil] (10.4,-51.5) -- (11.2,-38.7);
\draw[-{Stealth[length=1.8mm]}, pencil] (11.2,-36.3) -- (15.6,-21.5);
\draw[-{Stealth[length=1.8mm]}, pencil] (15.6,-19) -- (13.6,-7.5);
\draw[-{Stealth[length=1.8mm]}, pencil] plot[smooth, tension=0.55] coordinates {(18.5,-19.5) (24,-15.5) (29,-13.6)};
\draw[-{Stealth[length=1.8mm]}, pencil] (13,-37.5) -- (24.5,-32);
\draw[-{Stealth[length=1.8mm]}, pencil] (14.5,-69.5) -- (27.6,-63.5);
\draw[pencil] plot[smooth, tension=0.55] coordinates {(12.6,-42) (15,-50) (19,-57.5)};
\draw[-{Stealth[length=1.8mm]}, pencil] (27.6,-63.2) -- (29.5,-49);
\draw[-{Stealth[length=1.8mm]}, pencil] (30,-48) -- (37,-38.2);
\draw[-{Stealth[length=1.8mm]}, pencil] (33,-48.8) -- (43,-45);
\draw[-{Stealth[length=1.8mm]}, pencil] (37.2,-36.3) -- (38,-27.2);
\draw[pencil] (38.6,-37.8) -- (45.6,-31.5);
\node[pt] at (29.6,-81) {};
\draw[-{Stealth[length=1.8mm]}, pencil] (29.6,-81) -- (29.6,-66);
\draw[pencil] plot[smooth, tension=0.55] coordinates {(25,-72) (26.8,-77) (29.2,-80.5)};
\node[pt] at (51.2,-65) {};
\node[pt] at (75.5,-42) {};
\draw[-{Stealth[length=1.8mm]}, pencil] (51.2,-65) -- (53.6,-55.6);
\draw[-{Stealth[length=1.8mm]}, pencil] (55.2,-53) -- (61,-49.5);
\draw[-{Stealth[length=1.8mm]}, pencil] (62.6,-47.5) -- (68,-43.5);
\draw[-{Stealth[length=1.8mm]}, pencil] (51.2,-65) -- (60,-62);
\draw[-{Stealth[length=1.8mm]}, pencil] (61.5,-61) -- (68,-57.6);
\draw[-{Stealth[length=1.8mm]}, pencil] (69.5,-55.5) -- (73.2,-50.5);
\draw[-{Stealth[length=1.8mm]}, pencil] (74,-49.6) -- (75,-44.6);
\draw[pencil] (51.2,-65) -- (75.5,-42);
\draw[pencil] (52,-64) -- (75,-43.5);
\end{tikzpicture}
\end{document}
